\documentclass{article}
\usepackage{mathtools} 
\usepackage{fontspec}
\usepackage[UTF8]{ctex}
\usepackage{amsthm}
\usepackage{mdframed}
\usepackage{xcolor}
\usepackage{amssymb}
\usepackage{amsmath}
\usepackage{hyperref}


% 定义新的带灰色背景的说明环境 zremark
\newmdtheoremenv[
  backgroundcolor=gray!10,
  % 边框与背景一致，边框线会消失
  linecolor=gray!10
]{zremark}{注释}

% 通用矩阵命令: \flexmatrix{矩阵名}{元素符号}{行数}{列数}
\newcommand{\flexmatrix}[4]{
  \[
  #1 = \begin{pmatrix}
    #2_{11}     & #2_{12}     & \cdots & #2_{1#4}   \\
    #2_{21}     & #2_{22}     & \cdots & #2_{2#4}   \\
    \vdots      & \vdots      & \ddots & \vdots     \\
    #2_{#31}    & #2_{#32}    & \cdots & #2_{#3#4}
  \end{pmatrix}
  \]
}

% 简化版命令（默认矩阵名为A，元素符号为a）: \quickmatrix{行数}{列数}
\newcommand{\quickmatrix}[2]{\flexmatrix{A}{a}{#1}{#2}}


\begin{document}
\title{5.1注释}
\author{张志聪}
\maketitle

\begin{zremark}
  证明

  数域$K$上线性空间$V$内全体对称双线性函数与二次型函数的一一对应关系.
\end{zremark}
\textbf{证明：}
设对称双线性函数集合为$U$，二次型函数的集合为$V$，
定义函数$Q: U \to W$，接下来我们证明$Q$是双射即可。

\begin{itemize}
  \item 单射；

        对任意$f_1, f_2 \in U$，且$f_1 \neq f_2$，我们证明
        设$\epsilon_1, \epsilon_2, \cdots, \epsilon_n$是$V$的一组基，
        取$f_1, f_2$在这组基下的矩阵分别为$A = (a_{ij}), B = (b_{ij})$，$A \neq B$.

        对$X$的适当选择，可以得到矩阵中元素。
        对角线上的元素，取$\alpha = \epsilon_i$，
        于是
        \begin{align*}
          (e_i)^T A e_i = a_{ii}
        \end{align*}
        对于非对角线上的元素，$\alpha = \epsilon_i + \epsilon_j$，
        于是，由P340的结论
        \begin{align*}
          a_{ij}
           & = f_1(\epsilon_i, \epsilon_j)                                                                \\
           & = \frac{1}{2} [Q_{f_1}(\epsilon_i + \epsilon_j) - Q_{f_1}(\epsilon_i) - Q_{f_1}(\epsilon_j)]
        \end{align*}
        由于$A \neq B$，即存在$a_{ij} \neq b_{ij}$，
        即
        \begin{align*}
          a_{ij} \neq b_{ij} \\
          \frac{1}{2} [Q_{f_1}(\epsilon_i + \epsilon_j) - Q_{f_1}(\epsilon_i) - Q_{f_1}(\epsilon_j)]
           & \neq \frac{1}{2} [Q_{f_2}(\epsilon_i + \epsilon_j) - Q_{f_2}(\epsilon_i) - Q_{f_2}(\epsilon_j)]
        \end{align*}
        所以$Q_{f_1} \neq Q_{f_2}$，所以$Q$是单射.
  \item 满射。

        对任意$Q_f$，由P340的结论
        \begin{align*}
          f(\alpha, \beta) = \frac{1}{2} [Q_f(\alpha + \beta) - Q_f(\alpha) - Q_f(\beta)]
        \end{align*}
        可证明其是满射的。

\end{itemize}

\begin{zremark}
  $Q_f(\alpha) \equiv 0$则$f(\alpha, \beta) \equiv 0$的逆否命题：

  $f(\alpha, \beta) \not \equiv 0$，则$Q_f(\alpha) \not \equiv 0$。
\end{zremark}

\end{document}
